\section{Conclusion}

Completion of this laboratory work has provided insightful understanding of
preemptive multitasking with FreeRTOS on the ESP32 microcontroller and the design
principles behind real-time task synchronization.

This fourth laboratory work covered the port of the bare-metal cooperative application to
FreeRTOS preemptive tasks, extending it with synchronization mechanisms while preserving
the same three-layer architecture and task function structure. The key takeaways are:

\begin{enumerate}
  \item Porting a \textbf{cooperative bare-metal application to FreeRTOS} with minimal
    changes to task logic -- by encapsulating periodicity and offset entirely inside the
    scheduler service abstraction (\texttt{periodicTaskRunner} with \texttt{vTaskDelayUntil}),
    the task functions themselves remain identical in structure to the bare-metal version.
  \item Using \textbf{\texttt{vTaskDelayUntil}} for drift-corrected periodic execution,
    which guarantees that each task fires at precise intervals regardless of how long the
    task body takes to execute -- a significant improvement over relative \texttt{vTaskDelay}.
  \item Implementing \textbf{inter-task synchronization} with a binary semaphore for
    event-driven signaling (Task 1 notifies Task 2 on press detection) and a mutex for
    protecting shared statistical variables against concurrent access by Tasks 1, 2, and 3.
  \item Understanding the difference between a \textbf{binary semaphore} (used for
    signaling, no ownership) and a \textbf{mutex} (used for mutual exclusion, with
    ownership semantics and priority inheritance support in FreeRTOS).
  \item Structuring the application into \textbf{three distinct layers}: driver, service
    (scheduler + sync), and application -- maintaining the portability and modularity
    established in Laboratory Work 3.
  \item Gaining practical experience with \textbf{safe shared resource access} in a
    preemptive environment, where failing to protect global variables would lead to
    race conditions and corrupted statistics.
\end{enumerate}

This laboratory work builds directly on Laboratory Work 3, demonstrating how the same
application logic and architecture can be lifted from bare-metal to a full RTOS with
targeted, minimal changes -- a skill essential for developing scalable and maintainable
embedded firmware.
